We focus on ensemble datasets such as those produced by the Coupled Model Intercomparison Project (CMIP6)~\cite{eyring2016overview}, which contain climate projections generated under a variety of input assumptions. Each simulation produces a spatiotemporal time series, referred to as a model run. The following terminology is used throughout:
\begin{itemize}[noitemsep, topsep=0pt]
    \item \textbf{GCM} (Global Climate Model): Mathematical models that simulate the Earth's climate system, and are used to project future climate scenarios (e.g., MIROC6, ACCESS-CM2, etc.).
    \item \textbf{SSP} (Socioeconomic Pathways)~\cite{riahi2017shared}: Pathways that examine how global society, demographics, and economics might change over the next century. 
    In addition to the historical period, we consider three SSPs in this work: SSP245, SSP370, and SSP585, a low, medium high, and high greenhouse gas emissions future respectively. 
    \item \textbf{Model run}: A unique climate projection and spatiotemporal time series, representing a member of the climate ensemble generated by a specific combination (GCM, SSP). The full ensemble consists of multiple such runs.
\end{itemize}

\subsubsection{Design Goals}\label{ssec:design_requirements}
Throughout the project, in consultation with a climate scientist (paper co-author) specializing in regional climate variability and changes, we identified three key questions that drive exploratory analysis of ensemble climate datasets. These questions inform the system's design goals.
% Given a climate ensemble dataset, we list the motivating questions and the design requirements for \ClimateSOM.
\begin{tightItemize}
    \item[\req{R1}] {\normalfont \sffamily``Q: What is the general behavior of the set A of model runs?''}  \\ A climate scientist wishes to understand the general behavior of a single model run or a set of model runs and its general tendencies as it relates to the climate projection it generates. For example, do the model run(s) consistently produce projection of a certain quality or does it vary over time? \\ $\rightarrow$ \textbf{Requirement: }\textit{Provide a visual language and description to explore the behavior of a single set of model runs over the entire temporal domain.}
    \item[\req{R2}] {\normalfont \sffamily``Q: How does set A of model runs compare to set B of model runs?''} \\ A climate scientist wishes to compare the behavior of two sets of model runs A and B to understand the differences in their projections. For example, does set A of model runs produce similar projections compared to set B of model runs? What overlapping qualities do they share and what climate projection qualities are unique to run A or run B?\\ $\rightarrow$ 
    \textbf{Requirement: }\textit{Provide a visual language description and support to compare the behavior of two sets of model runs over the entire temporal domain.}
    \item[\req{R3}]{\normalfont \sffamily``Q: Which model runs or GCMs exhibit \textit{similar} behavior?''} \\ A climate scientist wishes to understand whether there are meaningful clusterings of ensemble runs and clustering of GCMs based on the SSP forcings. Do all ensemble runs produce similar results? Is the effect of SSP forcing similar on all GCMs?\\ $\rightarrow$ 
    \textbf{Requirement: }\textit{Support clustering of ensemble model runs based and clustering of GCM based on their SSP forcings over the entire temporal domain.} 
\end{tightItemize}
% The goal of \ClimateSOM is to provide answers to the above question while ensuring that they are readily decipherable.
% The goal of \ClimateSOM is to address the above questions with an interpretable visualization tool.
% \note{How to properly incorporate interpretablitity as a design requirement?}
% As outlined in~\cref{sec:intro} the main task

\subsubsection{Abstracting the Model Runs Into a Distribution}\label{ssec:abstraction}
Central to \ClimateSOM is a self-organizing map (SOM) trained to capture the spatial variability of a climate ensemble dataset, while retaining its interpretability.
%
An SOM is an unsupervised neural network that projects high-dimensional data into a 2D grid while preserving the topological structure, making it a useful tool dimensionality reduction (DR) of its input.
%
Training a SOM over a climate ensemble dataset yields a 2D SOM grid, each SOM node representing a spatial pattern over the studied region.
%
\ClimateSOM leverages this capability to represent each model run within the ensemble as a distribution.
%
Specifically, for each time step we first locate its best matching unit (BMU) in the SOM nodes, and allow the 2D position of that SOM node to represent that time step.
%
Aggregating these positions across all time steps produces a 2D distribution defined by the set of BMU's node locations.
%
\cref{fig:SOM} presents the construction of the distributional representation.

Noteworthy in this construction, however, is the projection step of the initial 2D SOM grid. 
%
While the SOM nodes maintain local smoothness, they can have distortions globally~\cite{sheridan_self-organizing_2011}; thus, defining the SOM node locations solely by their 2D grid location is not desirable.
%
To address this, we add a projection step of the SOM nodes to better preserve their relative structure and enhance its interpretability.
% 
%We thus aim to correct for these distortions by projecting the SOM nodes into a more interpretable orientation to better preserve the relative structure of the SOM nodes.
%
% Finally, we use the projected SOM node locations to construct the distributional representation as shown in~\cref{fig:SOM}.
\begin{figure}[H]
   \includegraphics[width=\linewidth]{SOM-10.pdf}
    \caption{The distributional abstraction for model runs in \ClimateSOM. Each model runs are abstracted to a distribution defined by a projected location of SOM Nodes, the \textit{SOM Node Space}.}
    \label{fig:SOM}
\end{figure}
\noindent
Formally, for an ensemble run with $n$ time steps, each a climate projection over some spatial domain, $R=[r_1, r_2, \ldots, r_n]$, we abstract this as a set of points in 2D space:
\begin{align*}
S=\{p_1,p_2,\ldots,p_n\}, \,\,\, p_i=\text{proj}(BMU(r_i))
\end{align*}
where $BMU(r_i)$ is the best matching SOM node for $r_i$, and proj$(\centerdot)$ maps the node to a 2D layout.
%
This distributional representation of the model runs serves as the backbone of \ClimateSOM.
%
% Therefore, enabling the interpretation of the projected space of the SOM nodes is principal to our workflow which we


% as a set of points $S=\{p_1, p_2, \ldots, p_n\}$ in an interpretable 2D space, where each $p_i$ is the dimensionality-reduced position of best matching unit (BMU) of $r_i$ within the trained SOM.
%
%Designed to achieve these qualities, the workflow of \ClimateSOM is shown in~\cref{fig:workflow} and is composed of the following steps:
\subsubsection{Workflow Overview}\label{ssec:workflow_overview}
The workflow contains one pre-step and three main steps (\cref{fig:workflow}).
\begin{itemize}[left=0pt, itemsep=0pt, labelsep=1em, align=left]
    \item[\step{S0} \textbf{Training the SOM}]
        The workflow begins with training the SOM on the dataset to ensure it has sufficient explanatory power---specifically, that the approximation via its BMU is adequate---and smoothness across the map.
        % A crucial assumption of \ClimateSOM is that SOM must capture \textit{sufficient} variance in the ensemble data, so that $r_i \approx BMU(r_i)$ is a reasonable approximation.
        % Careful consideration is required for the parameters for the SOM training and the metrics with which to evaluate and these are described in~\cref{ssec:som_training_parameters,ssec:som_training_metrics}.
        % The SOM is trained to capture the spatial variability of the input data set. The SOM is trained using the \textit{Kohonen} algorithm~\cite{kohonen1982self}, which is a type of unsupervised learning algorithm that is used to produce a low-dimensional representation of the input data set.
    \item[\step{S1} \textbf{Anchor}]
        Using the pseudo-2D smooth SOM nodes from Step 0, this step enables users to create the Adjusted SOM Node Space by projecting the SOM nodes into a 2D layout. To support interpretation and navigation, users can anchor specific nodes---fixing their positions while allowing others to adjust---resulting in a more interpretable spatial arrangement than the original SOM grid. 
        We refer to this refined space as the \textit{Adjusted SOM Node Space}.
        %
        % Here, instead of using common techniques like t-SNE\cite{van2008visualizing} or UMAP\cite{mcinnes2018umap}, we opt for Minimum-Distortion Embedding (MDE) \cite{agrawal2021minimum}, which generalizes t-SNE and UMAP and crucially allows the anchoring of points in the output 2D space.
        % %
        % When supplied with a anchoring constraint, MDE ensures that these points are \textit{fixed} in a given location in the output space, while reducing the distortion of the remaining points.
        % %
        % We prefer MDE over t-SNE and UMAP since often fail to preserve important characteristics of the embedded space like interpretable axes.
        %
        % We use MDE to allows users to iteratively construct the \textit{Adjusted SOM Node Space}, anchoring points within the embedded 2D space to ensure a more meaningful layout of the SOM nodes.
    \item[\step{S2} \textbf{Annotate}]
        Once the \textit{Adjusted SOM Node} Space is defined, users annotate it to produce the \textit{Annotated SOM Node Space}. An annotation consists of (1) a bounded region within the adjusted space and (2) a text label describing that region. These annotations serve as an abstraction layer to facilitate interpretation in the next step. \ClimateSOM incorporates LLMs to assist with this process.
    \item[\step{S3} \textbf{Analyze}]
        A given model run can now be described as a distribution within the \textit{Annotated SOM Node Space}, as well as percentages broken down by user-defined \textit{annotations} \req{R1}. Model runs can be compared either by directly comparing their distributions or by interpreting them as vector fields over the annotated space \req{R2}. Clustering can be applied to model runs based on these distributions, and to SSP forcings over each GCM by clustering changes in their respective distributions \req{R3}. Please see the details in~\cref{sec:workflow}.
\end{itemize}

\begin{figure*}[h]
    \centering
    \includegraphics[width=\textwidth]{workflow.pdf}
    \caption{The \ClimateSOM workflow. The workflow is composed of the SOM Training step and three main steps: Anchor, Annotate, and Analyze. In \step{S1}: Anchor, the user adjusts the SOM Node Space using iterative anchoring with Minimum-Distortion Embedding to produce the \textit{Adjusted SOM Node Space}(\subimage{2}). Dragged nodes are called \textit{anchors} and are fixed in place, while the rest of the nodes are adjusted. In \step{S2}: Annotate, the user annotates the SOM Node Space using manual inspection of the SOM Nodes or with the LLM Integration. Once annotated, in \step{S3} the user can
        use the \textit{Annotated SOM Node Space} (\subimage{3})
    as the basis to analyze the ensemble dataset.}
    \label{fig:workflow}
\end{figure*}



\noindent
Having outlined the workflow, we now address two important design choices: the use of SOMs and the intentional abstraction away from temporal ordering.

\vspace{3pt}
\noindent\textbf{SOM as the central tool.} 
We choose SOMs over other DR or clustering methods for the following two reasons:
\begin{enumerate}[noitemsep, topsep=0pt]
\item \textbf{Preservation of local topology}: SOMs preserve neighborhood relationships in lower-dimensions which enables meaningful spatial interpretations. Our distributional abstraction of model runs requires that the learned grid be continuous and neighboring nodes express similar spatial patterns. Other DR or clustering methods do not produce topologically smooth results in lower-dimensions and thus do not satisfy the conditions required in \ClimateSOM.
\item \textbf{Established use in climate science}: SOMs are widely used for analysis of climate data and are familiar to domain experts, aiding in its interpretation.
% \item \textbf{Visual continuity and smoothness}: Our data distributional abstraction of model runs requires that the learned grid must be continuous and neigh
\end{enumerate}

% Among the various possible methods to initially reduce the dimensionality of the input ensemble, \ClimateSOM opts for SOM for two main reasons: (1) SOMs are an established tool in climate science and are known to be successful in extracting salient features from spatiotemporal climate data and (2) \ClimateSOM relies on preserving the variability of the input ensemble via a pseudo-smooth 2D representation.  
%
% Maintaining smoothness in this representation is crucial, as discontinuities or irregularities can compromise the interpretability of the abstraction.
%
% While smoothness can be defined in various ways, we require that, locally at each node, the neighborhood maintains a consistent level of smoothness across the grid. 
%
% This ensures that the MDE embedding, powered by local neighborhood graphs, produces interpretable axes in 2D.
% \note{How do I properly define meaningful after MDE?}
%
% SOMs with the tuned parameters can guarantee this; however, other methods like hierarchical or \textit{k}-means clustering do not.

\vspace{3pt}
\noindent\textbf{Ignoring the temporal ordering.}
\ClimateSOM represents a model run $R$ as a set of points $S$, thus discarding temporal order (scrambling R yields the same S). This is an intentional trade-off for two reasons:
\begin{enumerate}[noitemsep, topsep=0pt]
    \item \textbf{Emphasis of broad patterns}: The goal is to analyze the general behavior of model runs over time rather than their specific sequence. This align with how GCM outputs are studied for broad differences in climate science.
    \item \textbf{Reducing complexity}: While preserving temporal order enables the use of time series analysis techniques, such as dynamic time warping, it also introduces complexity and amplifies the random variations inherent in climate ensemble data.
\end{enumerate}
This trade-off enables simpler interpretation while retaining the ability to extract broad patterns from model outputs.
% %
% (1) \
% %
% (2) while accounting for the temporal order allows for the application of time series  analysis and techniques like dynamic time warping, it is certain that it also introduces complexity and accentuates the random variations that climate ensemble data exhibits.
%
% Despite the loss of the temporal ordering, we argue that \ClimateSOM can still produce meaningful insights into the behavior of climate ensemble datasets.
